% Modernised reading — fonds Alexandre Grothendieck, shelfmark 49, pages 1-9.
%
% This is an INTERPRETATION, not a transcription. Everything the critical
% apparatus carried moves into footnotes. In any dispute of substance the
% transcription governs: whoever cites Grothendieck must cite batch-01.fr.tex.
%
% Pass: Opus 5 (claude-opus-5), 2026-08-30 — first pass, unchecked against the
% pages by a human. The folder was read whole; it holds one batch, pages 1-9,
% and page 1 is a cover.
%
% Notation fixed for the whole document. The manuscript's "simple sur S" is
% today's "lisse sur S", and that single word governs the whole first half:
% what page 2 says it remains to prove is smoothness. A' is his dual object;
% it is written A^vee here throughout, and the reading says once that the two
% are the same letter. The Picard functor is subscripted A/S everywhere,
% including at the three places (pages 2 and 4) where the pages write X/S or
% X_y for the same abelian scheme and its fibre. Nér.-Sév. is written NS.
%
% Substantive departures from the pages, each footnoted where it occurs:
% — page 3's Lemme is applied on the page to nothing: he states it, and the
%   sheets do not say to what. The reading supplies the application (G = O_{A^vee},
%   F = O_A, phi = phi_L), which is the only one under which the surrounding
%   argument closes, and footnotes that the supply is the edition's.
% — clause c) of that Lemme is partly illegible in its retained form; the
%   struck line above it is legible and states fibrewise injectivity. The
%   reading gives the struck form as the statement and says so.
% — page 2's chain dim A^vee_0 <= dim H^1 <= dim A_0 has two brace annotations,
%   the first of which is illegible except for the words "de Pic". The reading
%   supplies the justification (the tangent space of Pic at the origin) and
%   footnotes that the page does not say it.
% — page 6's theorem is stated for the kernel "exactly" the primitive classes;
%   the sheets prove one implication and name rigidity for the other on page 9,
%   where the folder stops after two lines. The reading says which half is on
%   the paper.
%
% What the folder announces and never establishes: the second route to the
% theorem proposed in the diagonal margin of page 2, which would avoid the
% torsion result — of it only the shape survives; the Lemme of page 7, whose
% statement is more than half illegible; and the end of the page-8 argument,
% cut off by the two lines of page 9. There is also no Corollaire 5: the
% numbering runs 1, 2, 3, 4, then 6.
\documentclass[11pt,a4paper]{article}
% Le préambule commun à toutes les transcriptions du fonds Grothendieck.
%
% Il tient en une idée : **l'incertitude se compose, elle ne se résout pas.**
% Une transcription qui lisse un mot douteux a détruit la seule chose pour
% laquelle on la faisait — un lecteur doit pouvoir savoir, à chaque endroit,
% ce qui était sur le feuillet et ce qui a été ajouté. Les six macros qui
% suivent sont donc l'appareil critique, et non de la décoration.
%
% Les mêmes commandes sont reconnues par `scripts/render.mjs`, qui produit la
% vue de lecture du site. Une macro ajoutée ici doit l'être aussi là-bas,
% faute de quoi le rendu échouera bruyamment — ce qui est voulu : un rendu
% silencieusement faux est pire qu'un rendu absent.

\ProvidesPackage{grothendieck}[2026/08/08 Transcriptions du fonds Grothendieck]

\usepackage{amsmath,amssymb,amsthm}
\usepackage{mathrsfs}
\usepackage{stmaryrd}
% Les diagrammes commutatifs sont partout dans ce fonds : c'est la langue
% même de la géométrie algébrique de Grothendieck, et une transcription qui
% les rendrait en prose ne transcrirait rien.
\usepackage{tikz-cd}
% Français par défaut — c'est la langue des feuillets — mais l'anglais est
% chargé aussi : la traduction et le résumé sont des documents anglais, et la
% typographie française (espaces avant les deux-points) y serait une faute.
% Ces documents basculent par \selectlanguage{english} en tête de corps.
\usepackage[english,french]{babel}
\usepackage{fontspec}
\usepackage{geometry}
\geometry{a4paper,margin=25mm,marginparwidth=28mm}
\usepackage{marginnote}
\usepackage{xcolor}
% ulem, not soul. `soul`'s \st and \ul are built on a character-by-character
% scan that XeLaTeX's font machinery breaks: the document fails with an
% undefined control sequence rather than degrading. ulem does the same two jobs
% and is Unicode-engine safe. `normalem` stops it hijacking \emph, which the
% transcriptions use for Grothendieck's own emphasis.
\usepackage[normalem]{ulem}
\usepackage{hyperref}
\hypersetup{colorlinks=true,linkcolor=black,urlcolor=black,pdfborder={0 0 0}}

% --- Métadonnées du lot -----------------------------------------------------
% Reprises telles quelles par le script de rendu, qui les affiche en tête de
% la vue de lecture. Elles fixent ce que le fichier prétend couvrir : sans
% elles, une transcription flotte sans point d'attache dans un fonds de seize
% mille feuillets.

\newcommand{\folder}[1]{\gdef\grothFolder{#1}}
\newcommand{\batch}[1]{\gdef\grothBatch{#1}}
\newcommand{\leaves}[2]{\gdef\grothFirst{#1}\gdef\grothLast{#2}}
\newcommand{\foldertitle}[1]{\gdef\grothFoldertitle{#1}}
\gdef\grothFolder{?}\gdef\grothBatch{?}\gdef\grothFirst{?}\gdef\grothLast{?}\gdef\grothFoldertitle{}

% --- Le résumé --------------------------------------------------------------
%
% Il ouvre la lecture modernisée, sous le titre « Résumé », et il est composé
% en petit. Ce n'est pas un ornement : c'est le seul passage du document qui
% ne soit pas des mathématiques, et un lecteur doit pouvoir le reconnaître
% avant de l'avoir lu — puis le sauter s'il connaît déjà le sujet, ou s'y
% arrêter s'il ne le connaît pas. Le filet qui le suit marque la reprise du
% corps.

\newenvironment{resume}
  {\small\setlength{\parskip}{0.45em}}
  {\par\medskip\hrule\medskip}

% --- Mots-clés ---------------------------------------------------------------
%
% En anglais, en clôture du résumé de la lecture modernisée : le vocabulaire
% moderne sous lequel chercher ce que les feuillets construisent. Le manifeste
% du site (`npm run manifest`) les extrait pour étiqueter la cote entière —
% cette ligne est la source unique des tags, il n'y a pas de fichier à côté.
\newcommand{\keywords}[1]{\par\smallskip\noindent
  {\footnotesize\textsc{Keywords} --- \textit{#1}}\par}

% --- L'appareil critique ----------------------------------------------------

% Le numéro de feuillet, en marge. C'est l'ancre qui permet de régler un
% désaccord sur une formule en regardant *un* feuillet plutôt qu'une chemise
% de six cents. Le site s'en sert aussi pour tourner les pages du fac-similé
% quand on fait défiler la transcription.
\newcommand{\leaf}[1]{%
  \marginnote{\footnotesize\color{gray!70}\textsf{#1}}%
  \label{leaf:#1}\ignorespaces}

% Illisible : on ne devine pas. Un mot inventé qui se lit comme les autres est
% le pire résultat possible.
\newcommand{\ill}{\textcolor{red!60!black}{[\,\ldots\,]}}

% Lecture douteuse : le mot est donné, et signalé comme incertain.
\newcommand{\uncertain}[1]{\uline{#1}}

% Ajout de l'éditeur — une lettre manquante, un mot sous-entendu.
\newcommand{\add}[1]{\textcolor{blue!50!black}{[#1]}}

% Biffé par Grothendieck lui-même. Ce qu'il a rayé fait partie du document :
% une rature dit souvent où la pensée a changé de direction.
\newcommand{\struck}[1]{\sout{#1}}

% Note du transcripteur, sur l'état matériel du feuillet.
\newcommand{\note}[1]{\par\smallskip{\small\color{gray!60!black}\textit{[#1]}}\par\smallskip}

% Note marginale de Grothendieck, distincte du corps du texte.
\newcommand{\marginal}[1]{\par\smallskip\noindent\hspace{1em}%
  {\small\color{gray!60!black}#1}\par\smallskip}

% --- Titre ------------------------------------------------------------------

\AtBeginDocument{%
  \begin{center}
    {\large\bfseries Fonds Alexandre Grothendieck --- cote n\textsuperscript{o}~\grothFolder}\\[2pt]
    {\itshape\grothFoldertitle}\\[4pt]
    {\small feuillets \grothFirst--\grothLast\ (lot \grothBatch)}\\[2pt]
    {\footnotesize\color{gray!60!black}Transcription établie d'après le fac-similé numérisé
     par l'Université de Montpellier.\\
     Les lectures incertaines sont soulignées, les illisibles notées [\,\ldots\,],
     les ajouts entre crochets.}
  \end{center}
  \vspace{1em}}

\endinput


\folder{49}
\pages{1}{9}
\dating{s.d.}
\watermark{Édition de démonstration\\interprétation personnelle de l'œuvre}
\foldertitle{Picard des schémas abéliens : notes manuscrites (s.d.) — lecture modernisée du dossier entier}

\begin{document}

\section*{Résumé}

\begin{resume}
Une famille d'objets géométriques, et la question de savoir si la famille de
leurs duals est encore une famille du même genre.

Une \emph{variété abélienne} est un objet géométrique dont les points
s'additionnent : la courbe elliptique en est le premier exemple, et en
dimension quelconque elles sont l'analogue algébrique du tore. À chacune est
attachée une seconde variété abélienne, sa \emph{duale}, de même dimension, et
la duale se construit à partir des fibrés en droites sur la première : c'est,
en gros, l'espace des manières de tordre une fonction sur la variété sans
rien changer à sa taille. Cet espace-là s'appelle le \emph{Picard} de la
variété.

Ce dossier ne travaille pas sur une variété mais sur une \emph{famille} : un
paramètre parcourt une base, et au-dessus de chaque valeur du paramètre il y a
une variété abélienne. On dit \emph{schéma abélien}. La question est simple à
énoncer et n'est pas facile : quand on prend la duale de chaque membre de la
famille, obtient-on encore une famille ? Il ne suffit pas que chaque duale
existe séparément. Il faut que la collection se recolle en un seul objet, et
surtout qu'elle varie \emph{sans à-coups} — que la duale ne saute pas ni ne
dégénère quand le paramètre bouge. C'est cette dernière exigence, la
lissité, qui fait tout le travail des quatre premiers feuillets.

L'idée qui arrive est de coincer une dimension entre deux inégalités. À un
fibré en droites $L$ assez positif sur la famille correspond une application de
la variété vers sa duale — translater $L$ et comparer au $L$ de départ. Cette
application a un noyau fini, donc elle ne perd pas de dimension : la duale est
au moins aussi grosse que la variété. En sens inverse, la duale ne peut pas
être plus grosse qu'un certain espace de cohomologie, et cet espace-là est
borné par la dimension de la variété — pour une raison purement algébrique, qui
est que la cohomologie d'un objet où les points s'additionnent est toujours
l'algèbre extérieure de son premier étage. Les deux inégalités se referment
l'une sur l'autre, l'égalité tombe partout, et avec elle la lissité.

Ce que le dossier établit ensuite est d'une autre nature, et c'est peut-être
le plus joli. Il y a une application naturelle qui envoie un fibré en droites
sur l'homomorphisme de la variété vers sa duale qu'il définit. Les fibrés que
cette application ne voit pas sont exactement ceux qui sont \emph{additifs} :
ceux dont la valeur en une somme de deux points est le produit des valeurs, au
sens précis où l'image réciproque par l'addition est le produit des deux
images réciproques par les projections. Autrement dit : le noyau d'une
application définie par la géométrie se décrit par une équation
fonctionnelle. C'est cette identité qui, aujourd'hui encore, sert de
définition à la partie neutre du Picard, et le dossier la démontre par une
descente — on remplace la base par la variété elle-même, ce qui rend
tautologique le point qu'on translatait.

Le dossier s'arrête en pleine démonstration, sur deux lignes qui invoquent la
rigidité. Il porte aussi, au feuillet 5, une phrase que l'on ne trouve pas
souvent sous cette plume : « J'ignore si le corollaire vaut sans supposer $A$
projectif sur $S$, et comme j'ignore si $\mathrm{Pic}_{A/S}$ existe dans ce
cas ». Toute l'hypothèse dont le dossier se sert est là, désignée par lui,
avec ce qu'il n'en sait pas.

\keywords{abelian scheme, dual abelian scheme, Picard scheme, Néron-Severi group scheme, polarisation, ample line bundle, isogeny, smoothness over a base, flatness criterion, cohomological flatness, exterior algebra, Hopf algebra theorem, primitive line bundle, theorem of the cube, rigidity lemma}
\end{resume}

\pagerange{1}{9}
\section*{Le dossier, sa charpente et ses notations}

Neuf feuillets, dont un de garde. Le feuillet 1 ne porte que le titre, de sa
main : \emph{Picard des schémas abéliens}. Les feuillets 2 à 9 sont une seule
suite, d'une seule écriture, et — fait rare dans ce fonds — aucun d'eux n'est
le verso d'autre chose.\footnote{L'inventaire de Montpellier avertit que « les
pages au verso des documents sont souvent sans rapport avec les notes
mathématiques ». Ce dossier-ci est celui où l'avertissement ne sert à rien : il
n'y a pas une ligne qui ne soit du sujet.}

L'argument tient en deux théorèmes et se lit en cinq stations :

\begin{itemize}
\item l'énoncé, et l'encadrement de la dimension par
  $\varphi_L$ d'un côté et par la cohomologie de l'autre (feuillet 2) ;
\item la fermeture de l'encadrement, et le lemme de platitude qui donne la
  lissité (feuillet 3) ;
\item l'algèbre extérieure de $R^1f_*\mathcal{O}_A$, qui est ce qui borne la
  cohomologie (feuillets 4 et 5) ;
\item les corollaires, et la seule phrase du dossier qui dise ce qu'il ne sait
  pas (feuillet 5) ;
\item le second théorème : l'accouplement
  $\mathrm{Pic} \times A \to A^{\vee}$, et la description de son noyau par
  l'identité d'additivité (feuillets 6 à 9).
\end{itemize}

Les notations sont fixées ici pour tout ce qui suit. $S$ est un schéma
localement noethérien, $A \to S$ un schéma abélien de dimension relative $n$,
$f : A \to S$ le morphisme structural et $\varepsilon$ sa section unité. On
écrit $A^{\vee}$ pour le schéma abélien dual, c'est-à-dire
$\mathrm{Pic}^{0}_{A/S}$ — c'est le $A'$ du manuscrit, et c'est aussi ce que le
dossier a précisément pour but de construire. On écrit $A_{0}$, $A^{\vee}_{0}$
pour les fibres en un point de $S$, $\mathrm{NS}_{A/S}$ pour le faisceau de
Néron--Severi relatif, et $m : A \times_S A \to A$ pour la loi de groupe.

Deux mots du manuscrit demandent d'être traduits, et l'un des deux commande
tout le premier théorème.

\emph{Simple} veut dire \emph{lisse}.\footnote{C'était le mot de l'époque —
EGA IV parlera encore de \emph{morphisme simple} avant que \emph{lisse} ne
s'impose. Ce n'est pas une nuance de style : quand le feuillet 2 écrit « il
reste à montrer qu'il est \emph{simple} sur $S$ », il désigne l'unique point
difficile du théorème, et un lecteur qui lirait \emph{simple} au sens
d'aujourd'hui ne verrait pas de quoi il s'agit.} Un morphisme lisse est une
famille sans à-coups : localement, elle ressemble à un produit. C'est
exactement ce qu'il faut établir pour que la collection des duales mérite le
nom de schéma abélien.

\emph{Nér.-Sév.} est le groupe de Néron--Severi, $\mathrm{Pic}/\mathrm{Pic}^{0}$
— les fibrés en droites comptés à déformation près. Sur une variété abélienne
il est sans torsion, et c'est de cela que le feuillet 2 se sert dès sa
troisième ligne.

Enfin, le manuscrit indice le foncteur de Picard tantôt par $A/S$, tantôt par
$X/S$, et écrit $X_{y}$ pour la fibre de $A$ ; il n'y a qu'un seul schéma
abélien dans le dossier et on écrit partout $A$.\footnote{La transcription
garde les $X$ là où ils sont, avec une note. Rien sur les feuillets ne suggère
une distinction : le $X$ est un reste de rédaction.}

\pagerange{2}{2}
\section*{Le théorème, et la première moitié de l'encadrement (feuillet 2)}

\textbf{Théorème.} \emph{Soit $A \to S$ un schéma abélien projectif. Alors}
$\mathrm{Pic}^{0}_{A/S}$ \emph{est un schéma abélien sur $S$, ouvert et fermé
dans} $\mathrm{Pic}_{A/S}$, \emph{de même dimension relative que $A$.}

Le manuscrit écrit d'abord $\mathrm{Pic}^{\tau}_{A/S}$, le biffe et écrit
$\mathrm{Pic}^{0}_{A/S}$. Les deux coïncident ici, et c'est la première chose
qu'il note : les fibres ont un groupe de Néron--Severi sans torsion, donc
$\mathrm{Pic}^{\tau}$, qui est par définition la partie de torsion modulo
$\mathrm{Pic}^{0}$, se réduit à $\mathrm{Pic}^{0}$.\footnote{La biffure est
donc une correction de notation et non de contenu, et elle enregistre le
moment où il s'aperçoit qu'il n'a pas besoin des deux objets. C'est aussi ce
qui rend le corollaire 3 du feuillet 5 — $\mathrm{Pic}^{\tau}$ ouvert et fermé
— non redondant : cet énoncé-là vaut plus largement.}

De cela découlent, sans travail, la propreté sur $S$ et la fermeture dans
$\mathrm{Pic}_{A/S}$ ; l'ouverture, elle, est générale. Ce qui reste, et qui
est tout, est la \emph{lissité} sur $S$ — les fibres étant déjà connues
géométriquement connexes.

Voici l'instrument. On choisit $L$ ample relativement à $S$, ce qui donne une
section $s$ de $\mathrm{Pic}_{A/S}$ et l'homomorphisme
\[
  \varphi_{L} : A \longrightarrow A^{\vee}, \qquad
  \varphi_{L}(t) = \left[\, \tau_{t}^{*}L \otimes L^{-1} \,\right],
\]
où $\tau_{t}$ est la translation par $t$.\footnote{Le manuscrit écrit
$\tau_{t}^{*}(L)\,L^{-1}$, le groupe de Picard en notation additive ; on garde
le produit tensoriel, qui est l'usage. Il note aussi que $\varphi_{L}$ envoie
section unité sur section unité, donc est un homomorphisme de schémas en
groupes — c'est le lemme de rigidité, qu'il invoque par un simple « (cf.) » et
qui reviendra explicitement au feuillet 9.} C'est ce que l'on appelle
aujourd'hui la \emph{polarisation} associée à $L$.

Sur une fibre, $\varphi_{L}$ est surjectif. La raison qu'il donne est que son
noyau est fini, par le critère d'amplitude de Weil ; un homomorphisme de
variétés abéliennes à noyau fini est surjectif sur sa cible dès que les
dimensions le permettent, d'où
\[
  \dim A_{0} \;\leqslant\; \dim A^{\vee}_{0}.
\]
Reste à borner l'autre côté. Il écrit la chaîne
\[
  \dim A^{\vee}_{0} \;\leqslant\; \dim H^{1}(A_{0}, \mathcal{O}_{A_{0}})
  \;\leqslant\; \dim A_{0},
\]
et annote chacune des deux inégalités d'une accolade.

La première tient à ce que l'espace tangent à l'origine de
$\mathrm{Pic}$ est $H^{1}(A_{0}, \mathcal{O}_{A_{0}})$ : déformer un fibré en
droites au premier ordre, c'est se donner une classe dans ce groupe, et une
variété est de dimension au plus celle de son espace
tangent.\footnote{L'annotation sous cette accolade est illisible sauf les deux
derniers mots, « de Pic ». La justification donnée ici est celle sous laquelle
l'inégalité est vraie ; le feuillet ne l'écrit pas.}

La seconde est la clé du dossier, et c'est un fait d'algèbre pure. Sur un
objet où les points s'additionnent, la cohomologie du faisceau structural est
une \emph{algèbre de Hopf} graduée : la loi de groupe lui donne un coproduit.
Le théorème de structure de ces algèbres force
\[
  H^{*}(A_{0}, \mathcal{O}_{A_{0}}) \;=\;
  \Lambda^{*} H^{1}(A_{0}, \mathcal{O}_{A_{0}}),
\]
donc $H^{i}$ est nul dès que $i$ dépasse $\dim H^{1}$ ; et comme $H^{i}$ est
aussi nul au-delà de $\dim A_{0}$ pour des raisons de dimension, on obtient
$\dim H^{1} \leqslant \dim A_{0}$.\footnote{C'est bien ce que porte le
feuillet : sous la seconde accolade, deux lignes qui se chevauchent, où l'on
lit « th.\ de Hopf » au-dessus de la nullité de $H^{i}$ pour $i > \dim A_{0}$.
Le second mot de la ligne inférieure est douteux ; l'argument, lui, ne l'est
pas, car les deux ingrédients nommés ne se combinent que d'une seule façon.}

Le bas du feuillet porte, en diagonale sur toute la marge gauche et sur une
vingtaine de lignes, une seconde route vers le même théorème, qui se passerait
du résultat de torsion invoqué en tête. Elle passe par la platitude de
$\varphi_{L}$ et par le fait que $\mathrm{Pic}^{0}$ est ouvert et fermé dans
$\mathrm{Pic}$. Il n'en reste que la forme : l'écriture y est plus rapide que
partout ailleurs dans le dossier et la liaison est perdue.\footnote{On ne
reconstruit pas cette route ici. Le peu qui se lit — « Si on ne veut pas se
servir du résultat de la torsion », puis « surjectif, et de plus \emph{plat} »,
puis « c'est donc un schéma abélien » — dit l'intention et l'ordre, pas les
étapes.}

\pagerange{3}{3}
\section*{La fermeture de l'encadrement, et la platitude (feuillet 3)}

Les deux inégalités se referment : $\dim A_{0} \leqslant \dim A^{\vee}_{0}$ et
$\dim A^{\vee}_{0} \leqslant \dim A_{0}$ donnent l'égalité partout. Il en tire
trois choses d'un coup : $A^{\vee}_{0}$ est lisse, sa dimension est celle de
$A_{0}$, et
\[
  \dim H^{1}(A_{0}, \mathcal{O}_{A_{0}}) \;=\; n,
\]
donc toute la structure de $H^{*}(A_{0}, \mathcal{O}_{A_{0}})$ est
connue.\footnote{« d'où la structure totale de $H^{*}$ » : puisque
l'algèbre est extérieure sur son premier étage, connaître la dimension de
celui-ci les détermine toutes, $\binom{n}{q}$ en degré $q$.}

Reste à passer des fibres à la famille. La lissité de $A^{\vee}$ sur $S$
demande la platitude, et c'est pour l'obtenir qu'il énonce un lemme. Sa forme
retenue est en partie illisible ; sa forme biffée, juste au-dessus, ne l'est
pas, et c'est elle qu'on donne.\footnote{Le feuillet porte, biffé d'un trait :
« $G_{s} \to \varphi_{s*}(F_{s})$ est injectif ». La clause c) qui le remplace
commence par « tout élément de $G_{s}$ dont l'image inverse en les
$x \in \varphi_{s}^{-1}(y)$ s'annule », et la fin manque. Les deux disent la
même chose, la seconde ponctuellement ; on retient la première, qui est
complète.}

\textbf{Lemme.} \emph{Soient $X$ et $Y$ de type fini sur $S$ localement
noethérien, $\varphi : X \to Y$ un $S$-morphisme, $F$ et $G$ quasi-cohérents
sur $X$ et $Y$, et $G \to \varphi_{*}F$ un morphisme. Soient $s \in S$ et
$y \in Y$ au-dessus de $s$. On suppose : $F$ plat sur $S$ ; $G$ cohérent ; et
$G_{s} \to \varphi_{s*}(F_{s})$ injectif au-dessus de $y$. Alors $G$ est plat
sur $S$ en $y$.}

C'est un critère de platitude par injection dans un objet plat : un faisceau
cohérent qui se plonge fibre à fibre dans un plat est plat. Le manuscrit ne dit
pas à quoi il l'applique — il l'énonce et passe au feuillet suivant. La seule
application sous laquelle l'argument des feuillets 2 et 3 se referme est
celle-ci : on prend $X = A$, $Y = A^{\vee}$, $\varphi = \varphi_{L}$,
$F = \mathcal{O}_{A}$ et $G = \mathcal{O}_{A^{\vee}}$.\footnote{C'est la
lecture de l'édition, et le feuillet ne l'écrit pas. Ce qui la rend
pratiquement forcée est la ligne qui précède immédiatement le lemme :
« pour tout $x' \in A^{\vee}_{0}$ il existe $x \in A_{0}$ au-dessus et un
homomorphisme $\mathcal{O}_{A^{\vee}_{0}, x'} \to \mathcal{O}_{A_{0}, x}$
\emph{injectif} ». C'est exactement l'hypothèse d'injectivité fibre à fibre du
lemme, écrite sur les anneaux locaux. Le qualificatif de cet homomorphisme est
douteux sur la page ; ce que l'énoncé demande est qu'il soit local et
injectif.} L'hypothèse de platitude porte alors sur $\mathcal{O}_{A}$, qui est
plat sur $S$ parce que $A$ est un schéma abélien, et la conclusion est que
$\mathcal{O}_{A^{\vee}}$ l'est aussi.

Plat, de présentation finie, à fibres lisses : $A^{\vee}$ est lisse sur $S$.
Le théorème est acquis.

\pagerange{4}{5}
\section*{L'algèbre extérieure en famille (feuillets 4 et 5)}

Le feuillet 4 porte ce que le manuscrit appelle, en anglais dans le texte,
\textbf{Cor.\ Main 1}, et qui est l'énoncé dont le second maillon de
l'encadrement était le cas particulier sur une fibre.

\textbf{Corollaire.} \emph{Soit $A \to S$ propre et plat, $f$ le morphisme de
projection. Alors le morphisme d'algèbres}
\[
  \Lambda^{*}\!\left(R^{1}f_{*}\mathcal{O}_{A}\right)
  \;\xrightarrow{\ \sim\ }\; R^{*}f_{*}\mathcal{O}_{A}
\]
\emph{est un isomorphisme, et $R^{1}f_{*}\mathcal{O}_{A}$ est localement libre
sur $S$, de rang égal à la dimension relative de $A$ sur $S$.}

La démonstration a trois temps, et l'ordre en est instructif.

D'abord, le cas d'un corps est acquis — c'est le théorème de Hopf du feuillet
2. Ensuite, la section unité $\varepsilon$ donne
\[
  R^{1}f_{*}\mathcal{O}_{A} \;=\;
  \mathcal{H}om\!\left(\Omega^{1}_{\varepsilon}, \mathcal{O}_{S}\right),
\]
donc un module localement libre, et — c'est la formulation qu'il retient —
$\mathcal{O}_{A}$ est \emph{cohomologiquement plat} pour $f$ en dimension
$1$.\footnote{Il écrit « coh.\ plat », abréviation qu'il ne développe nulle part
dans le dossier. Le membre de droite est d'abord écrit
$\varepsilon^{*}(\Omega^{1}_{A/S})$, puis biffé au profit de
$\Omega^{1}_{\varepsilon}$ : les deux désignent le même conormal le long de la
section unité.} La platitude cohomologique est ce qui permet de comparer la
cohomologie de la famille à celle de chaque fibre.

Enfin le carré, qui est l'argument proprement dit :

\begin{tikzcd}[column sep=small, row sep=small, nodes={font=\scriptsize}]
  \Lambda^{*}\!\left(R^{1}f_{*}\mathcal{O}_{A}\right)_{y}
    \arrow[r, "\alpha"] \arrow[d, "\text{surj.}"']
  & \left(R^{*}f_{*}\mathcal{O}_{A}\right)_{y} \arrow[d] \\
  \Lambda^{*}H^{1}(A_{y}, \mathcal{O}_{A_{y}})
    \arrow[r, "\simeq"']
  & H^{*}(A_{y}, \mathcal{O}_{A_{y}})
\end{tikzcd}

La ligne du bas est l'isomorphisme de Hopf sur la fibre. La colonne de gauche
est surjective par construction. La colonne de droite est le morphisme de
changement de base, et la platitude cohomologique le rend bijectif. Il ne
reste alors à $\alpha$ qu'à être surjectif fibre à fibre, ce qui, pour un
morphisme de modules gradués de même rang au-dessus d'un anneau local, suffit
à en faire un isomorphisme. Le feuillet 5 le conclut d'une ligne : « c'est un
isomorphisme, cqfd ».\footnote{Le raccord entre les deux feuillets est
elliptique — la phrase du bas du feuillet 4 est en partie illisible, et le
haut du feuillet 5 reprend au milieu d'une justification. Le pas manquant est
celui indiqué ici : surjectivité sur les fibres plus égalité des rangs.}

Trois corollaires suivent immédiatement, et une remarque qui vaut mieux
qu'eux.

\textbf{Corollaire 2.} \emph{Si $L$ est très ample relativement à $S$, alors
$\varphi_{L} : A \to A^{\vee}$ est surjectif et fidèlement plat, et son noyau
est un groupe fini et plat sur $S$.} C'est dire que $\varphi_{L}$ est une
\emph{isogénie}, et que $\ker \varphi_{L}$ — le groupe qu'on note aujourd'hui
$K(L)$ — est fini plat.\footnote{Le manuscrit prend soin d'écrire « très ample
(et non seulement ample) », et souligne la parenthèse. L'énoncé vaut en fait
pour $L$ relativement ample ; la restriction qu'il s'impose ici est celle sous
laquelle sa démonstration passe, non celle que l'énoncé demande.}

\textbf{Corollaire 3.} \emph{$\mathrm{Pic}^{\tau}_{A/S}$ est ouvert et
\emph{fermé}.} Il ajoute, entre crochets, la seule chose qui compte :
l'ouverture vaut toujours, la fermeture est propre aux schémas
abéliens.\footnote{L'exposant est ici tracé comme une croix, alors que le
$\tau$ des feuillets 2 et 6 est net ; c'est le même. L'énoncé n'est pas
redondant avec le théorème, puisqu'il porte sur $\mathrm{Pic}^{\tau}$ sans
supposer que $\mathrm{NS}$ soit sans torsion.}

\textbf{Corollaire 4.} \emph{Existence de $\mathrm{NS}_{A/S}$.} Il l'écrit
ainsi, en trois mots, et c'est en effet immédiat : un quotient par un
sous-groupe ouvert et fermé existe. C'est le corollaire 3 qui le porte.

\pagerange{5}{5}
\section*{Ce qu'il ignore (feuillet 5)}

Entre la fin de la démonstration et les corollaires, une remarque :

\begin{quote}
J'ignore si le corollaire vaut sans supposer $A$ projectif sur $S$, et comme
j'ignore si $\mathrm{Pic}_{A/S}$ existe dans ce cas\ldots
\end{quote}

La phrase s'achève sur deux mots que la page ne soutient pas. Ce qu'elle dit
est net : toute la construction repose sur la projectivité, qui sert deux
fois — à disposer d'un $L$ ample, et à savoir que le foncteur de Picard est
représentable.

L'hypothèse n'est pas anodine, et c'est ce qui donne son prix à la remarque.
Un schéma abélien n'est pas nécessairement projectif sur sa base ; et la
construction du dual d'un schéma abélien quelconque, sans cette hypothèse,
est aujourd'hui acquise.\footnote{On énonce ici deux faits, et non une
généalogie : l'édition ne prétend pas dire quand il les a sus ni par qui, et
ces feuillets ne sont pas datés. Ce qui est du dossier, et seulement cela,
c'est la question — posée à l'endroit exact où l'hypothèse est utilisée.}

\pagerange{6}{7}
\section*{L'accouplement, et les classes additives (feuillets 6 et 7)}

Le feuillet 6 ouvre le second mouvement, et il change de sujet : il ne s'agit
plus de construire $A^{\vee}$ mais de comprendre l'application qui va de tous
les fibrés en droites vers les homomorphismes.

Elle s'écrit de deux façons, qu'il numérote :
\[
  (*) \qquad \mathrm{Pic}_{A/S} \longrightarrow
  \mathcal{H}om_{S\text{-gr}}(A, A^{\vee}),
\]
\[
  (**) \qquad \mathrm{Pic}_{A/S} \times_{S} A \longrightarrow A^{\vee},
  \qquad (\xi, t) \longmapsto \tau_{t}^{*}\xi - \xi .
\]
La seconde est la première « curryfiée » : au lieu d'envoyer $\xi$ sur une
application, on envoie le couple sur une valeur.\footnote{Il porte en marge,
sur deux lignes reliées par une accolade, « hom.\ ponctuel » et « hom.\ de
groupes » : c'est exactement la distinction entre les deux formes, et la
flèche va de la seconde vers la première.}

\textbf{Théorème.} \emph{$(*)$ est un homomorphisme de $S$-groupes — c'est-à-dire
que $(**)$ est un accouplement — et son noyau est exactement l'ensemble des
sections $\xi$ vérifiant}
\[
  m^{*}\xi \;=\; \mathrm{pr}_{1}^{*}\xi + \mathrm{pr}_{2}^{*}\xi .
\]

C'est le cœur de la seconde moitié. L'identité de droite dit que $\xi$ est
\emph{additive} au sens le plus littéral : sa valeur en une somme est la somme
de ses valeurs. En notation multiplicative,
$m^{*}L \cong \mathrm{pr}_{1}^{*}L \otimes \mathrm{pr}_{2}^{*}L$ ; c'est ce
qu'on appelle aujourd'hui un fibré \emph{primitif}, et c'est la
caractérisation usuelle de $\mathrm{Pic}^{0}$.\footnote{Le manuscrit écrit
$\pi$ pour la loi de groupe et note à côté du cadre « $\pi : A \times_{S} A \to
A$ \ldots\ la loi de composition \ldots\ i.e.\ $\sigma$ » — il hésite entre deux
lettres pour le même morphisme. On écrit $m$, qui est l'usage. Le point
d'interrogation qu'il porte après « homomorphisme de $S$-groupes » est de sa
main et porte sur cet énoncé-là.}

Le corollaire qu'il en tire est celui pour lequel tout ce qui précède a été
fait :
\textbf{Corollaire 6.}\quad
\[
  \mathrm{NS}_{A/S} \longrightarrow \mathcal{H}om_{S\text{-gr}}(A, A^{\vee}),
\]
compatible avec les structures de groupe. Puisque le noyau de $(*)$ est
$\mathrm{Pic}^{0}$, l'application descend au quotient et y devient
injective : le groupe de Néron--Severi se plonge dans les homomorphismes de
$A$ vers son dual.\footnote{Le dossier s'arrête à cette injection. On sait
aujourd'hui en décrire l'image — ce sont les homomorphismes symétriques, du
moins sur un corps algébriquement clos. Le feuillet ne l'aborde pas, et il n'y
a pas de \emph{Corollaire 5} : la numérotation du dossier passe de 4 à 6.}

Le feuillet 7 est la démonstration, et il est le plus abîmé du dossier. Ce qui
s'en lit tient en une observation : l'objet construit est trivial sur
$A^{\vee} \times \varepsilon$ et sur $\varepsilon \times A$, donc trivial sur
le produit, qui est encore un schéma abélien. Il énonce ensuite un
\textbf{Lemme} dont plus de la moitié est illisible, portant sur une
trivialisation au-dessus d'un corps algébriquement clos, et le renvoie à plus
bas : « Cf.\ \ldots\ et sera prouvé plus bas ».\footnote{Ce lemme n'est prouvé
nulle part dans le dossier. Le renvoi est vers des pages qui n'y sont pas, ou
qui n'ont pas été écrites.}

\pagerange{8}{9}
\section*{La descente à la base tautologique (feuillets 8 et 9)}

Le feuillet 8 démontre l'équivalence annoncée, dans un sens, et par le procédé
qui rend l'énoncé presque tautologique.

Le point de départ est
\[
  \varphi_{\xi} = 0
  \quad \Longleftrightarrow \quad
  m^{*}\xi - \mathrm{pr}_{1}^{*}\xi - \mathrm{pr}_{2}^{*}\xi = 0,
\]
lue sur le diagramme

\begin{tikzcd}[column sep=small, row sep=small, nodes={font=\scriptsize}]
  A & A \times_{S} A \arrow[l, "\mathrm{pr}_{2}"'] \arrow[r, "m"]
    \arrow[d, "\mathrm{pr}_{1}"] & A \\
  & A &
\end{tikzcd}

L'idée est de changer de base pour le paramètre lui-même. On pose $S_{1} = A$,
$A_{1} = A \times_{S} S_{1}$, et l'on dispose alors sur $A_{1}$ d'une
\emph{section tautologique} $t_{1}$ : au-dessus du point $t$ de la base, elle
vaut $t$. Dire que $\varphi_{\xi}$ est nul revient exactement à dire que cette
section-là ne bouge pas $\xi$ :
\[
  t_{1}^{*}(\xi_{1}) = \xi_{1}, \qquad \xi_{1} = \xi \times_{S} S_{1}.
\]
Le quantificateur « pour tout $t$ » a disparu dans la base ; c'est tout le
bénéfice de la manœuvre.\footnote{Le feuillet écrit d'abord cette égalité,
puis biffe une ligne qui la répétait avec la section unité. La descente est
son argument, et il ne la commente pas.}

Le calcul qui suit vérifie que la combinaison
$m^{*}L - \mathrm{pr}_{1}^{*}L - \mathrm{pr}_{2}^{*}L$ est triviale le long de
la section unité. En restreignant par $i_{1}$,
\[
  i_{1}^{*}\!\left(m^{*}L - \mathrm{pr}_{1}^{*}L - \mathrm{pr}_{2}^{*}L\right)
  = (m\,i_{1})^{*}L - (\mathrm{pr}_{1} i_{1})^{*}L
  - (\mathrm{pr}_{2} i_{1})^{*}L \;=\; 0,
\]
les deux premières composées étant l'identité et la troisième la section unité
$\varepsilon_{A/S}$ : le premier terme et le deuxième s'annulent l'un l'autre,
et le troisième est trivial.\footnote{Le manuscrit porte ces trois
identifications sous trois accolades, « id », « id » et
« $\varepsilon_{A/S}$ ». Ce sont elles qui annulent le membre de droite, et
c'est la seule chose que le feuillet démontre complètement dans cette
section.}

Ce que ce calcul donne est une moitié de l'énoncé : si $\varphi_{\xi} = 0$,
l'obstruction à l'additivité est triviale le long des deux sections. Pour
conclure qu'elle est triviale tout court — et donc pour l'autre moitié de
l'équivalence — il faut un argument de rigidité, et c'est là que le dossier
s'arrête. Le feuillet 9 ne porte que deux lignes, où l'on lit « cette
rigidité pour la \ldots », et le reste de la page est blanc.\footnote{Un fibré
en droites sur un produit, trivial sur chacun des deux facteurs, est trivial :
c'est ce que les traités appellent le théorème du cube, dont le lemme de
rigidité est l'ingrédient. Le mot qu'il écrit est bien \emph{rigidité} ; ce
qu'il en aurait fait, la page ne le dit pas, et l'édition ne le complète pas
à sa place.}

\end{document}
